\bookchapter{第十八章：データバックアップと移行}

\chapterintro{一、まず確認：スマホ内のどれがなくなると本当に消えるのか}

\booksection{異なるデータの「生存状態」は異なり、区別する必要があります：}

スマホローカルにあり、なくなると本当に消えるもの

\begin{itemize}
  \item \textbf{：写真動画（クラウド同期していない場合）、メモ、SMS、通話履歴、ダウンロードしたファイル。}Appサーバーにあり、機種変更して再ログインすればあるもの
  \item \textbf{：WeChatの友達リスト、Taobaoの注文、銀行Appの取引記録。これらはスマホ内ではなく対応する企業のサーバーに保存されています。}グレーゾーン、特別な操作が必要なもの
  \item \textbf{：WeChatチャット履歴（WeChatサーバーに保存されていますが、自分でサーバーからダウンロードできません）、スマホバンキングのデバイス証明、2段階認証App内のシークレット（Google Authenticatorなど）。}だから重点的に保護すべきは、第一類と第三類です。
\end{itemize}

二、全体バックアップ：一度にすべて対応

\booksection{全体バックアップはスマホの完全な「スナップショット」を撮るようなもの——現在のシステム設定、デスクトップレイアウト、Appデータ、写真、連絡先、すべてパッケージ化。}

方案一：パソコンローカルバックアップ（最も信頼できる、パソコン保有者に推奨）

\booksubsection{iPhone}

\begin{itemize}
  \item \textbf{iPhone}：ケーブルでスマホをパソコンに接続し、iTunes（旧Windows版）またはFinder（Mac）を開いて、デバイスページで「今すぐバックアップ」を選択。「バックアップを暗号化」にチェックしないと、パスワードやヘルスケアデータはバックアップされません。
  \item \textbf{Android}パソコンバックアップの利点はクラウド容量を買う必要がなく、復元速度がクラウド復元よりはるかに速い。欠点は定期的にケーブル接続してバックアップする必要がある——ほとんどの人ができない。
  \item 方案二：クラウドバックアップ（最も手間いらず、すべての人に推奨）
\end{itemize}

\booksubsection{iPhone}

\begin{itemize}
  \item \textbf{iPhone}：設定 ユーザの名前 iCloud iCloudバックアップ オン。電源接続·Wi-Fi接続·ロック画面の状態で自動バックアップされます。無料の5GBではまず足りないので、有料の50GBプランへのアップグレードがおすすめ。
  \item \textbf{Android}クラウドバックアップの最大の利点は「気にしなくていい」、自動で完了。最大の欠点は復元時にダウンロードに長時間かかり、無料容量は通常足りない。
  \item 三、分類バックアップ戦略：重要なものは別途保存
\end{itemize}

\booksection{全体バックアップは最後の砧ですが、一部のものは重要すぎて別途保存する価値があります：}

写真と動画（最も貴重なデジタル資産）

\booksubsection{二重保険方案}

\begin{itemize}
  \item \textbf{：iCloud写真/Googleフォトの自動同期をオンにしつつ、数ヶ月おきに手動で写真をパソコンや外付けHDDにコピー。クラウド1部、ローカル1部、相互にバックアップ。}お金を使いたくない方案
  \item \textbf{：Googleフォトは無料で15GB提供、「圧縮画質」に設定して容量節約可能。または1ヶ月おきにケーブルで写真をすべてパソコンにコピー。}連絡先（見落としがちだが極めて重要）
\end{itemize}

\booksubsection{iPhone連絡先はデフォルトでiCloudと同期、同期をオンにしていれば追加バックアップ不要。ただしより安全にするなら、iCloud.comでvCardファイルを書き出してパソコンに保存。}

\begin{itemize}
  \item Android連絡先は通常Googleまたはブランドクラウドと同期。連絡先Appの設定で.vcfファイルを書き出せば、完全なバックアップになります。
  \item WeChatチャット履歴（個別処理、最重要）
\end{itemize}

\booksubsection{これは個別に説明する必要があります、全体バックアップではWeChatチャット履歴は多くの場合不完全で、機種変更ツールでは全く移動できないから。}

方法一：新スマホに移行

\begin{itemize}
  \item \textbf{。WeChat マイ 設定 一般 チャット履歴の移行とバックアップ 移行 スマホ/タブレットに移行。新旧スマホを同じWi-Fiに接続し、QRコードをスキャンして転送。注意：移行完了まで旧スマホのWeChatをログアウトしない。}方法二：パソコンにバックアップ
  \item \textbf{。パソコンにWeChat PC版をインストール、左下の三本線 移行とバックアップ バックアップと復元 チャット履歴をパソコンにバックアップ。機種変更後に新スマホに復元可能。}方法三：クラウドにバックアップ
  \item \textbf{。WeChatはiCloudやスマホクラウドスペースへのバックアップにも対応しますが、注意：このバックアップは現在のスマホでのみ復元可能、クロスデバイス復元不可。機種変更する場合は、前2つの方法を使う必要があります。}四、機種変更ガイド：旧スマホから新スマホへ
\end{itemize}

\booksection{新スマホへの変更は嬉しいことですが、データが移行できないと頭痛に。3つのケースをそれぞれ対応：}

同ブランドアップグレード（最も簡単）

\booksubsection{iPhone 新しいiPhone}

\begin{itemize}
  \item \textbf{：起動アクティベーション時に旧スマホを横に置くと、「クイックスタート」が自動ポップアップ、旧スマホで新スマホのパターンをスキャンするだけ。すべてのデータ、設定、Appレイアウトがそのまま移行、Wi-Fiパスワードも自動的に引き継がれます。約20〜30分。}Android同ブランド
  \item \textbf{：ブランド標準の機種変更ツールを使用。Huaweiは「スマホクローン」、Xiaomiは「Xiaomi機種変更」、OPPOは「スマホ引越し」、vivoは「相互転送」。手順はほぼ同じ：新旧スマホに同じAppをインストール、QRコードで接続、移行したいものを選択、転送開始。}クロスブランド機種変更（Android間、異ブランド）
\end{itemize}

\booksubsection{サードパーティの機種変更ツールまたは旧ブランド標準の機種変更Appを使用（ほとんどのブランドの機種変更Appは他のAndroidスマホへの移行にも対応）。移行できるもの：連絡先、写真、SMS、通話履歴、一部App。移行できないもの：App内のログイン状態（再ログイン必要）、WeChatチャット履歴（WeChat独自の移行機能を使用）。}

クロスプラットフォーム機種変更（AndroidiPhone、またはその逆）

\booksubsection{Android iPhone}

\begin{itemize}
  \item \textbf{：Appleの「iOSに移行」Appがあり、Androidスマホでダウンロード後、新iPhoneアクティベーション時に「Androidデバイスからデータを移行」を選択。連絡先、写真、SMSなどの基本データは移行できるが、Appデータは移行不可。WeChatチャット履歴はWeChat独自の移行機能を使用。}iPhone Android
  \item \textbf{：各Androidブランドに対応する移行ツールがありますが、正直クロスプラットフォーム移行体験は同プラットフォームほど良くない。最も重要なことを先に個別に：連絡先書き出し、写真バックアップ、WeChatチャット履歴移行。残りはゆっくり。}五、バックアップの黄金の法則：3-2-1原則
\end{itemize}

\booksection{これはデータセキュリティ分野の古典的原則で、シンプルで覚えやすい：}

3

\begin{itemize}
  \item \textbf{：重要なデータは少なくとも3部保持（原本1部、バックアップ2部）。}2
  \item \textbf{：バックアップは2種類の異なるストレージメディアを使用（例：クラウド1部、外付けHDD1部）。}1
  \item \textbf{：少なくとも1部のバックアップを異なる物理的位置に配置（火災、盗難などで同時にすべてのデータが失われるのを防止）。}一般人にとっては、「クラウド1部、パソコン1部」で既に十分。複雑にする必要はなく、本当に実行することが重要。
\end{itemize}

WeChatチャット履歴の特別な注意

\begin{tcbwarn}{WeChatチャット履歴はどのスマホ機種変更ツールでも移行できません！全体バックアップのカバー範囲外です。WeChat標準の「チャット履歴移行」機能を使用するか、パソコン版WeChatにバックアップする必要があります。多くの人が新スマホに機種変更した後にチャット履歴がすべてなくなったことに気づき、その時はもう手遅れ——旧スマホは既に出荷時リセットされている可能性。}
練習：完全なバックアップ訓練を実施

\end{tcbwarn}

\begin{tcbpractice}{iCloudクラウドバックアップまたはブランドクラウドバックアップが有効か確認、前回のバックアップ時間を確認。}
\begin{enumerate}
  \item 手動で写真をバックアップ：ケーブルでパソコンに接続、DCIMとPicturesフォルダをパソコンにコピー。
  \item 連絡先をバックアップ：iPhoneユーザーはiCloud.comでvCardを書き出し。Androidユーザーは連絡先Appの設定で.vcfファイルを書き出し。
  \item WeChatチャット履歴をバックアップ：WeChatを開く マイ 設定 一般 チャット履歴の移行とバックアップ パソコンにバックアップ。
  \item 書き出した連絡先vcfと写真フォルダを外付けHDDやUSBメモリにコピーし、「3-2-1」の「1部を別の場所にオフラインで」を完了。
  \item 第十九章：トラブルシューティングと自己救助
\end{enumerate}

\end{tcbpractice}
